\chapter{Crisis en el espacio real}\label{sec:crisis-en-el-espacio-real}

La atmósfera a bordo del buque insignia terrestre era muy diferente a la de mi visita anterior. El comandante me dejó en una sala de reuniones con dos guardias armados mientras hablaba con la Inteligencia Terrana sobre la traición de Ula. Afirmó que era para mi protección, pero a juzgar por su comportamiento, esa no era exactamente la verdad. El dúo me miró con ojos sospechosos todo el tiempo, sin esbozar una sonrisa ni decir una palabra.

Afortunadamente, Rykov no tardó mucho. No pude evitar notar que su rostro estaba desencajado por la preocupación cuando regresó. Fuera lo que fuese lo que acababa de descubrir, tenía la sensación de que no eran buenas noticias.

—Ya volviste. ¿Te dijeron qué van a hacer respecto al sabotaje? —pregunté.

El comandante suspiró. —Al parecer, la Tierra tiene problemas más grandes que resolver.

Mis antenas se movieron confundidas. —¿Qué podría ser más grande que una traición? ¡Tienen que hacer algo con la presidenta!

—Están de acuerdo en que hay que hacer algo, pero la tienen en la mira por otras razones —respondió—. Ha estado alentando actos de violencia contra nuestros civiles. Ha habido suficientes incidentes en los últimos días como para que la Tierra cierre sus puertos espaciales. No podemos arriesgarnos a un ataque terrorista en nuestro territorio.

 —¿Ataque terrorista? —supuse que la frase debía ser alguna jerga militar terrestre con la que no estaba familiarizado—. ¿Qué significa eso?

Los ojos del comandante Rykov se abrieron de par en par. —No... Eh, bueno, es una forma de violencia contra los civiles. Eventos con víctimas masivas que se planifican, se ejecutan públicamente y tienen como objetivo aterrorizar a un determinado grupo de personas.

El hecho de que los terranos tuvieran un término designado para tales ataques implicaba que ocurrían con cierta regularidad. Me estremecí ante la idea de civiles asesinados a plena luz del día con el único propósito de la crueldad. A pesar de todas las guerras que los jatari libramos en nuestros primeros años, la violencia nunca fue tan insensata. Sin embargo, el comandante hablaba de estos «ataques terroristas» como si fueran algo que simplemente sucedía, como un desastre natural.

 —¿Por qué me mira así? —Rykov frunció el ceño y se cruzó de brazos—. No intente actuar como si esas cosas nunca ocurrieran en la Federación.

Me moví con incomodidad. —Bueno, es que no es así.

—Yo no lo veo de esa forma. ¿Conoce nuestra embajada en la capital de la Federación? —hizo una pausa, esperando mi asentimiento—. Los manifestantes la asaltaron y tomaron a los diplomáticos como rehenes. Ahora están atrincherados dentro, amenazando con volar el lugar si intentamos algo.

 —¿Qué? ¿Por qué harían eso? Los humanos han sido muy amables con nosotros, no hay razón para derramar sangre.

—Intente decirles eso. Tal como yo lo veo, son terroristas, pero ese sigue siendo un tema de debate en casa. Deberíamos recuperar la embajada, no intentar negociar con esa gente. ¡Ni siquiera le hablan a nadie que sea humano!

Me di cuenta de que había algo que Rykov no me decía. Se preocupaba por la gente bajo su mando y por la preservación de la vida en general, pero nunca hasta el punto de que se notara tanto en su voz. Si esta misión era personal para él, necesitaba saber por qué. Era un comandante competente, pero nada nubla el juicio tanto como las emociones.

—No quiero entrometerme, pero... —le puse una mano en el hombro y noté que los guardias se tensaban ante el contacto—. ¿Conoce a alguien allí?

 —S-sí, mi hermano vive en la embajada. Trabaja para el Departamento de Estado —respondió Rykov—. ¿Era tan obvio?

Fruncí el ceño. —Quizás podamos ayudar de alguna manera. ¿Cuáles son sus órdenes?

—Bueno, primero teníamos que regresar a la Tierra con los refugiados. Pedí unirme al equipo táctico desplegado fuera de la embajada, pero me dijeron, y cito textualmente: «Está demasiado involucrado en esto». ¿Los peces gordos realmente esperaban que me quedara al margen?

—Las órdenes son órdenes. Tendremos que pensar en un plan cuando lleguemos.

—Voy un paso por delante de usted, general. Les dije que usted se había ofrecido a ayudarme y que me había pedido que lo escoltara hasta la embajada. Sería un insulto para los jatari si me negaba. Y, ¿sabe?, me concedieron el permiso.

 —¿Les mintió a sus superiores?

—La vida de mi hermano está en juego. No mentí, solo exageré la verdad.

—Esa es la definición de mentir.

—No es el punto. De todos modos, el buque insignia partirá hacia la capital inmediatamente y las otras naves llevarán a los refugiados de regreso a la Tierra. No lo obligaré a ir, pero agradecería su ayuda. ¿Se unirá a nosotros?

No había duda de que debía acompañar al comandante. Iba más allá de deberle la vida o considerarlo un amigo. Tras la traición de la presidenta y el abandono de la Federación, los humanos necesitaban ver que todavía tenían aliados. Necesitaban saber que no estaban solos.

—Por supuesto. Es lo mínimo que puedo hacer —respondí.

Rykov sonrió. —Excelente. Vayamos al puente y salgamos de aquí.

Los guardias estuvieron a mi lado en un instante, empujándome hacia adelante como si fuera ganado. Mi paciencia para el trato de prisionero se estaba agotando; estaba deseoso de ayudar a los humanos, pero no a expensas de mi dignidad.

Hice lo mejor que pude para imitar la pose de brazos cruzados de Rykov, mirando fijamente al guardia de la derecha. —No me toque. Estoy ocupado.

La figura corpulenta del hombre era intimidante y no tenía duda de que podría molerme a golpes en una pelea. Sabía que mi reprimenda solo podría conducir a un trato más duro, pero lo miré a los ojos de todos modos. Esperaba una mirada irritada, pero en cambio vi el atisbo de una sonrisa burlona en sus labios.

—Consiguió hacer sonreír a Mac, general. Ni siquiera yo puedo hacer eso —rio el comandante.

Mac resopló. —Eso es porque usted no tiene gracia.

—¡Oye! Tenga cuidado —el tono de Rykov era ligero y juguetón, así que dudé que su advertencia fuera seria—. Pero el general tiene razón. Es perfectamente capaz de caminar por su cuenta.

Mac asintió y dio un paso atrás. Su compañero hizo lo mismo y se quedó unos pasos detrás de nosotros. Mis pensamientos vagaron mientras pasábamos por los estrechos pasillos hacia el puente. Deseaba que el comandante hubiera despedido a los guardias por completo, pero me conformaría con tenerlos fuera de mi espacio personal. Si los humanos todavía no confiaban en mí, ese era su problema.

—Parece molesto. Lo siento, en realidad no es por usted —dijo Rykov—. Ahora tengo seguridad las veinticuatro horas. Un grupo extremista antihumano puso una recompensa de un millón de créditos por mi cabeza.

Mi paso se tambaleó cuando la conmoción me atravesó las venas. —¡¿Qué?! ¿Por qué por ti?

Suspiró. —Porque fui yo quien lanzó la nanobomba de nanitos.

Entramos en un silencio cómodo. Mi mente todavía estaba dando vueltas por la última revelación. Si hubiera una recompensa de un millón de créditos por mí, no confiaría ni en que mi propia tripulación no me liquidara. Para que el comandante deambulara por la nave, debía tener una fe inquebrantable en la lealtad de sus hombres.

En mi última visita al buque insignia, el puente se veía envuelto en un caos de actividad y charlas, al menos antes de que Rykov los llamara al orden. Pero hoy los miembros de la tripulación ya esperaban en sus puestos y el rumbo ya estaba trazado en la computadora. Al ver las expresiones sombrías a mi alrededor, temí que los últimos días hubieran minado el espíritu de los humanos.

—Muy bien, parece que estamos listos. El curso hacia la capital ya está ingresado en el sistema —Rykov asintió con satisfacción—. ¡Activen el motor de distorsión!

Mi estómago dio un vuelco cuando la nave se deslizó hacia el hiperespacio; incluso después de años de experiencia en viajes superlumínicos, las náuseas nunca desaparecieron. Recuerdos de soldados recién reclutados de la Federación vomitando hasta las tripas flotaron en mi mente y me reí para mis adentros. El comandante levantó la vista de su holocuadro, probablemente confundido por encontrar algo divertido en un momento así.

Antes de que pudiera preguntar, una voz severa se escuchó en la frecuencia de emergencia: —Regresen de inmediato. Las naves terrestres no autorizadas tienen prohibido el ingreso al espacio de la Federación hasta nuevo aviso.

Rykov frunció el ceño y presionó algunos botones en su pantalla. —¿Desde cuándo? Saben que somos miembros de la Federación, así que no veo cómo se nos puede negar el acceso.

—La Tierra suspendió los vuelos interplanetarios esta mañana y el Senado aprobó una restricción temporal a los visitantes terrestres como respuesta —fue la réplica—. Por lo que sabemos, cerraron sus puertos espaciales porque están planeando un ataque.

—¡Eso es ridículo! —protestó Rykov—. Mire, estamos en una misión de rescate. Vamos a atracar y usted no se va a interponer en nuestro camino.

—Están invadiendo este sistema. Si entran al espacio real, les dispararemos.

—¿Qué es eso? Se están entrecortando, no puedo oírlos —el comandante finalizó la transmisión, sacudiendo la cabeza con disgusto—. Idiotas. Levanten los escudos a máxima potencia y entren al espacio real.

Sentí que la nave se estremecía bajo mis pies al salir del hiperespacio y mi expresión se transformó en una de horror. —¡No, debe regresar! Los sistemas de defensa planetaria los derribarán.

—Me gustaría verlos intentarlo —gruñó Rykov.

Bueno, esto era simplemente maravilloso. Parecía que la presidenta Ula podría estar cumpliendo su deseo de una guerra con los humanos después de todo.